\section{Alimentos funcionales}

\definition{Alimento funcional}{Es un alimento que incluye componentes más allá de los nutrientes básicos y ayuda a la salud.}

\subsection{Pro, pre y postbióticos (y más)}

Prebióticos: Sirven como alimento para la microbiota intestinal.
Probióticos: Se componen de organismos vivos como tal, los cuales se añaden a la microbiota intestinal y ayudan en el trabajo de digestión.
Simbiótico: Es un alimento que contiene tanto prebióticos como probióticos
Postbiótico: Es un organismo no vivo que aporta antígenos al sistema inmune.

\section{Cumplimiento en etiquetado sobre probióticos}

Cuando en un producto se declara que este contiene probióticos, se debe demostrar que tal microorganismo debe ser capaz, no solo de atravesar el proceso digestivo sin morir, sino además de que debe de sobrevivir una población lo suficientemente grande para colonizar el área objetivo.
El proceso digestivo consta de:
1. Masticar
	1. Se aplica acción mecánica
	2. Actua la enzima amilasa salival
		1. Trabaja a pH neutro
		2. Rompe enlaces glucocídicos alfa ($\alpha$) como los de la glucosa, no beta ($\beta$) como los de la celulosa
	3. El alimento se empieza a considerar bolo alimenticio
2. Deglutir
	1. Se pasa por la laringe al esófago
	2. Se realizan movimientos peristálticos y sinuoides para emujar el bolo alimenticio al estómago
3. Estómago
	1. Aquí actuan varios agentes
		1. Ácido gástrico ($HCl$)
			1. pH 1.5-2
		2. Proteasas (Pepsinas)
		3. Lipasas
		4. Sales biliares
			1. Actuan sobre lípidos
4. Intestinos
	1. Aquí actuan otra serie de enzimas
	2. Entorno neutro-alcalino